
\documentclass{article}
\usepackage{enumitem}


\usepackage{graphicx,comment,framed}
\usepackage{roundbox}
\usepackage{fancybox}
\usepackage{tikz}
\usepackage{color}
\usepackage[hidelinks]{hyperref}
\usepackage{framed}
\usepackage{amsthm,amssymb,amsmath}
%\usepackage[colorlinks,linkcolor=blue,citecolor=blue]{hyperref}
\definecolor{shadecolor}{cmyk}{0,0,0,0}
\usepackage{listings}
\usepackage{xepersian}
\usepackage[noend]{algpseudocode}


%--------------------- page settings ----------------------

\settextfont[Scale=1.1]{XB Niloofar}
\setdigitfont[Scale=1.1]{XB Niloofar}
\defpersianfont\sayeh[Scale=1.1]{XB Niloofar}
\addtolength{\textheight}{3.2cm}
\addtolength{\topmargin}{-22mm}
\addtolength{\textwidth}{3cm}
\addtolength{\oddsidemargin}{-1.5cm}


%------------------------ Environments ------------------------------------

\newtheorem{قضیه}{قضیه}
\newtheorem{لم}{لم}
\newtheorem{مشاهده}{مشاهده}
\newtheorem{تعریف}{تعریف}


%-------------------------- Notations ------------------------------------

\newcommand{\IR}{\ensuremath{\mathbb{R}}} 
\newcommand{\IZ}{\ensuremath{\mathbb{Z}}} 
\newcommand{\IN}{\ensuremath{\mathbb{N}}} 
\newcommand{\IS}{\ensuremath{\mathbb{S}}} 
\newcommand{\IC}{\ensuremath{\mathbb{C}}} 
\newcommand{\IB}{\ensuremath{\mathbb{B}}} 

\newcommand{\bR}{\mathbb{R}}
\newcommand{\cB}{\mathcal{B}}
\newcommand{\cO}{\mathcal{O}}
\newcommand{\cG}{\mathcal{G}}
\newcommand{\rM}{\mathrm{M}}
\newcommand{\rC}{\mathrm{C}}
\newcommand{\rV}{\mathrm{V}}

\newcommand{\lee}{\leqslant}
\newcommand{\gee}{\geqslant}
\newcommand{\ceil}[1]{{\left\lceil{#1}\right\rceil}}
\newcommand{\floor}[1]{{\left\lfloor{#1}\right\rfloor}}
\newcommand{\prob}[1]{{\mbox{\tt Pr}[#1]}}
\newcommand{\set}[1]{{\{ #1 \}}}
\newcommand{\seq}[1]{{\left< #1 \right>}}
\newcommand{\provided}{\,|\,}
\newcommand{\poly}{\mbox{\rm poly}}
\newcommand{\polylog}{\mbox{\rm \scriptsize polylog}\,}
\newcommand{\comb}[2] {\left(\!\!\begin{array}{c}{#1}\\{#2}\end{array}\!\!\right)}




\newcounter{probcnt}
\newcommand{\مسئله}[1]{\stepcounter{probcnt}\bigskip\bigskip{
 	\large \bf مسئله‌ی \arabic{probcnt}$\mbox{\bf{.}}$ \ #1} \bigskip}

\newcommand{\fqed}[1]{\leavevmode\unskip\nobreak\quad\hspace*{\fill}{\ensuremath{#1}}}

\newenvironment{اثبات}
	{\begin{trivlist}\item[\hskip\labelsep{\em اثبات.}]}
	{\fqed{\square}\end{trivlist}}

\newenvironment{حل}
	{\begin{trivlist}\item[\hskip\labelsep{\bf حل.}]}
	{\fqed{\blacktriangleright}\end{trivlist}}

\ifdefined\hidesols
	\newsavebox{\trashcan} % uncomment the following line to hide solutions
	\renewenvironment{حل}{\begin{latin}\begin{lrbox}{\trashcan}}{\end{lrbox}\end{latin}}
\fi


%------------------------- Header -----------------------------

%------------------------- Header -----------------------------

\newcommand{\سربرگ}[3]{
	\parindent=0em
	\begin{shaded}
		
		\rightline{ 
			\makebox[6em][c]{
				\includegraphics[height=1.6cm]{aut.png}
		}}
		\vspace{-.2em}
		{\scriptsize\bf دانشکده‌ی مهندسی کامپیوتر}
		\hfill {\small
			مدرس: دکتر امیرمزلقانی  \ 
		}\\[-5em]
		\leftline{\hfill\Large\bf 
			کاربرد های جبر خطی
		}\\[.7em]
		\leftline{\hfill\bf 
			نیم‌سال دوم ۹۶-۹۷
		}\\[1.7em]
		\hrule height .12em
		
		\normalsize
		\vspace{1mm} #1
		\hfill \small  #3
		\vspace{1mm} 
		\hrule height .1em
		
		\vspace{-0.5em} 
		\hfill {\sayeh\large #2} \hfill
	\end{shaded}
	\begin{large}

توجه:
\end{large}
\\
\\
\begin{itemize}
\item
این تمرین از مباحث مربوط به فصل5 ،6و7 طراحی شده است که شامل 15سوال اجباری و 4 سوال امتیازی است که نمره سوال های امتیازی فقط به نمرات تمرین شما کمک می کند.
\item 
در برخی سوالات چند قسمتی از شما خواسته شده است به تعدادی از قسمت ها به اختیار پاسخ دهید ،پاسخگویی به بیش از مقدار تعیین شده نمره اضافی ندارد ،از بقیه قسمت ها می توانید برای تمرین بیشتر استفاده کنید.
\item 
در پایان نیز تعدادی سوال برای تمرین بیشتر در نظر گرفته شده است ،که به حل آن ها نمره ای تعلق نمی گیرد صرفا برای تمرین توسط خود شما در نظر گرفته شده است. 
\item 
اگه سوالی داشتین از طریق
\begin{latin}
\href{mailto:aut.la2018@gmail.com}{\[aut.la2018@gmail.com\]}
\end{latin}
 حتما بپرسید.

\item پاسخ های تمرین را در قالب یک فایل به صورت الگوی زیر آپلود کنید.
\begin{latin}
9531000\_steve\_McManaman\_HW4.pdf
\end{latin}
\item \color{red}
مهلت تحویل جمعه  1تیر 1397 ساعت 23:54:59
\end{itemize}

}




%\settextfont{Yas}
\begin{document}
\سربرگ{تمرین فصل5،6و7}{}{}
\clearpage
\مسئله{} یکی از ماتریس های زیر  را به اختیار انتخاب کنید ابتدا چند جمله ای سرشت نما را برای آن بیابید سپس مقدار ویژه و بردار های ویژه را برای آن  مشخص کنید  در نهایت صورت قطری شدن آن را قطری کنید.

$$\begin{bmatrix}
-1&4&-2\\
-3&4&0\\
-3&1&3
\end{bmatrix} \qquad
\begin{bmatrix}
2&2&-1\\
1&3&-1\\
-1&-2&2
\end{bmatrix}
 $$
 
 \مسئله{}
5 مورد از گزاره های زیر را به اختیار ثابت کنید:
\begin{enumerate}
	\item 
	نشان دهید اگر 
	$\lambda$
	مقدار ویژه ماتریس واون پذیر 
	$A$
	باشد آنگاه 
	$\lambda^{-1}$
	مقدار ویژه ماتریس 
	$A^{-1}$
	است.
	\item 
	نشان دهید اگر 
	$A^2=0$
	انگاه تنها  مقدار ویژه 
	$A$
	صفر است.
	\item 
		$\lambda$
		مقدار ویژه از 
		$A$
		است اگر و فقط اگر 
		مقدار ویژه ای از 
		$A^{T}$
		باشد.
		\item
		
		نشان دهید 
		$A$
		و 
		$A^{T}$
		جند جمله ای سرشت نمای مشابه ای دارند.
		\item 
		با توجه به الگوریتم 
		$QR$
		ثابت گنید اگر 
		$A=QR$
		باشد که 
		$Q$
		معکوس پذیر است آنگاه 
		$A$
		با 
		$A_1=RQ$
		متشابه است.
		\item 
		فرض کنید 
		$A$
		یک ماتریس حقیقی 
		$n\times n$
		باشد که 
		$A^T=A$.
		نشان دهید اگر برای 
		$x$
		های غیر صفری در 
		$\mathbb{C}^n$،
			$Ax=\lambda x$
		باشد آنگاه 
		$\lambda$
		حقیقی است و در واقع قسمت حقیقی 
		$x$
		بردار ویژه 
		$A$
		است.
		\item 
		برای بردار های 
		$u,v$
		در 
		$\mathbb{R}^n$
		ثابت کنید:
		$$\parallel u+v\parallel ^2 +\parallel u-v \parallel=2\parallel u\parallel^2+2\parallel v\parallel^2$$
		\item
		اگر 
		$U,V$
		دو ماتریس 
		$n\times n$
		متعامد باشند،نشان دهید 
		$UV$
		نیز یک ماتریس متعامد است.
\end{enumerate}


\مسئله{}
فرض کنید 
$A$
ماتریس 
$n\times n$
 باشد که مجموع درایه های تمام سطر های آن
$s$
باشد ثابت کنید 
$s$
مقدار ویژه ای از 
$A$
است.

	
	\مسئله{}
	({\bf سوال امتیازی })
	برای هر اسکالر 
	$a,b,c$
	نشان دهید:
	$$A=\begin{bmatrix}
	b&c&a\\
	c&a&b\\
	a&b&c
	
	\end{bmatrix},
	B=\begin{bmatrix}
	c&a&b\\
	a&b&c\\
	b&c&a
	
	\end{bmatrix}
	,C=
	\begin{bmatrix}
	a&b&c\\
	b&c&a\\
	c&a&b
	\end{bmatrix}
	$$
	همگی متشابهند و اگر
	$BC=CB$
	باشند آنگاه 
	$A$
	دو مقدار ویژه صفر دارد.



\مسئله{}({\bf سوال امتیازی })
اگر 
$$A=
\begin{bmatrix}
-1&1&1&-1\\
1&-1&-1&1\\
1&-1&-1&1\\
-1&1&1&-1
\end{bmatrix}$$
باشد،
$A^2,A^6$
را محاسبه کنید.


\مسئله{}فرض کنید
$\varepsilon=\{e_1,e_2,e_3\}$
پایه استاندارد برای 
$\mathbb{R}^3$
و 
$\mathcal{B}=\{b_1,b_2,b_3\}$
پایه ای برای فضای برداری 
$V$
باشد و 
$T:\mathbb{R}^3\longrightarrow V$
یک تبدیل خطی باشد که:
$$T(x_1,x_2,x_3)=(x_3-x_1)b_1-(x_1+x_2)b_2+(x_1-x_2)b_3$$
\begin{enumerate}
	\item 
	$T(e_1),T(e_2)$
	و 
	$T(e_3)$
	را محاسبه کنید.
	\item 
	$[T(e_1)]_{\mathcal{B}},[T(e_2)]_{\mathcal{B}}$
	و 
	$[T(e_3)]_{\mathcal{B}}$
		را محاسبه کنید.
		\item 
		ماتریس تبدیل 
		$T$
		را تحت پایه های 
		$\varepsilon,\mathcal{B}$
		بیابید.
		

 
	\end{enumerate}


\مسئله{}ثابت کنید مجموع درایه های روی قطر اصلی هر ماتریس قطری شدنی برابر است با مجموع مقادیر ویژه آن ماتریس.



\مسئله{} یک دیگر از روش هایی  زمانی که تقریبی از بردار ویژه در دسترس باشد می شود با آن مقادیر ویژه را یافت 
روش خارج قسمت ریلی 
({\lr rayleigh quotient })
است. 

مشاهده کردیم اگر 
$Ax=\lambda x$
آنگاه 
$x^TAx=x^T(\lambda x)=\lambda(xTx)$
و در این صورت خارج قسمت ریلی 
$$R(x)=\frac{x^TAx}{x^Tx}$$
ّبرابر
$\lambda$
خواهد بود.اگر 
$x$
به حد کافی به به یک بردار ویژه 
$\lambda$
نزدیک باشد آنگاه این خارج قسمت به 
$\lambda$
نزدیک خواهد شد.زمانی که 
$A$
متقارن باشد خارج قشمت ریلی 
$R(x_k)=(x^T_kAx_k)/(x^T_kx_k)$
با دقتی دو برابر نسبت 
$\mu_k$
در روش توانی عمل خواهد کرد این موضوع را برای ماتریس و بردار اولیه زیر نشان دهید:
$$A=\begin{bmatrix}
5&2\\
2&2
\end{bmatrix},x_0=\begin{bmatrix}
1\\
0
\end{bmatrix}$$

\مسئله{}({\bf سوال امتیازی })
اگر 
$W$
یک زیر فضا از فضای ضرب داخلی 
$V$
باشد،
$u\in V$
را در نظر بگیرید نشان دهید
$v\in W$
تصویری از 
$u$
بر روی 
$W$
است به طوری که   
 $$u=v+v'\qquad for\  \ some\ \ v' \in W^{\perp}$$
 اگر و فقط اگر 
 $$\parallel u-v\parallel \leq \parallel u-w\parallel\quad,\quad for \ \ every \ \ w\in W $$

 
\مسئله{}({\bf سوال امتیازی }) فرض کنید 
$W_1,W_2$
زیر فضایی از فضای ضرب داخلی 
$V$
باشد آنگاه نشان دهید:
\begin{enumerate}
	\item $(W_1+W_2)^{\perp}=W_1^{\perp}\cap W_2^{\perp}$
	\item $(W_1\cap W_2)^{\perp}=W_1^{\perp}+W_2^{\perp}$
\end{enumerate}



 \مسئله{} فرض کنید 
 $y=\begin{bmatrix}
 7\\
 9
 \end{bmatrix}$
 و 
 $u_1=\begin{bmatrix}
 \frac{1}{\sqrt{10}}\\
 -\frac{3}{\sqrt{10}}
 \end{bmatrix}$
 و 
 $W=span\{u_1\}$.
 
 
 $proj_W \boldsymbol{y},(UU^T)y$
  را حساب کنید.
 
 \مسئله {} فرض کنید 
 $W$
 زیر فضایی از 
 $\mathbb{R}^n$
با پایه متعامد 
$\{w_1,w_2,\cdots,w_p\}$
و همچنین فرض کنید 
$\{v_1,v_2,\cdots,v_q\}$
پایه ای متعامد برای 
$W^{\perp}$
باشد.
\begin{enumerate}
\item 
چرا 
$\{w_1,w_2,\cdots,w_p,v_1,v_2,\cdots,v_q\}$
یک پایه متعامد است ؟
\item 
چرا 
$span$
مجموعه قسمت 1
$\mathbb{R}^n$
را تولید می کند؟
\item 
نشان دهید
$dim W+dim W^{\perp}=n$.

	
\end{enumerate}
 \مسئله {}یکی از ماتریس های زیر را به اختیار انتخاب و برای فضایی ستونی آن یک پایه متعامد پیدا کنید:
 $$\begin{bmatrix}
 	3&-5&1\\
 	1&1&1\\
 	-1&5&-2\\
 	3&-7&8
 \end{bmatrix} \qquad
 \begin{bmatrix}
 	1&2&-1\\
 	-1&1&-4\\
 	-1&4&-3\\
 	1&-4&7\\
 	1&2&1
 \end{bmatrix}
 $$
 
 
\مسئله{} تمام جواب های کوچکترین مربعات را برای تساوی 
$Ax=b$
بیابید.
$$A= \begin{bmatrix}
1&1&0\\
1&1&0\\
1&0&1\\
1&0&1
\end{bmatrix},
b=\begin{bmatrix}
1\\
3\\
8\\
2
\end{bmatrix}
$$
\مسئله{}فرض کنید 
$A$
یک ماتریس 
$m\times n$
باشد که ستون هایش مستقل خطی هستند و
 $b\in \mathbb{R}^n$.
 
 با استفاده  از
 روش نرمال یک فرمول برای 
$\hat{b}$
که تصویر 
$b$
بر روی 
$Col \ \ A$
هست بیابید.

\مسئله{}نشان دهید اگر 
$A$
یک ماتریس
$n\times n$
مثبت معین باشد،آنگاه یک ماتریس مثبت معین 
$n\times n$
مانند 
$B$
وجود دارد که 
$A=BB^T$.


\مسئله{} ماتریس 
$A$
  را در نظر بگیرید،
$\lambda_1$
بزرگترین مقدار ویژه آن و
$u_1$
بردار ویژه یکه متناظر با 
$\lambda_1$
است،ثابت کنید بزرگترین مقدار 
$x^TAx$
با توجه به قیود:
$$x^Tx=0\qquad x^Tu_1=0$$
 برابر 
 $\lambda_2$
 است که 
 $\lambda_2$
 دومین مقدار ویژه بزرگ 
 $A$
 است. همچنین این بزرگترین مقدار زمانی اتفاق می افتد که 
 $x$
 برابر 
 $u_2$
 که بردار ویژه یکه متناظر با 
 $\lambda_2$
 است،باشد.

\مسئله{} تجزیه 
$SVD$
ماتریس زیر را به دست آورید.(راهنمایی:
ماتریس
$\begin{bmatrix}
-\frac{1}{3}&\frac{2}{3}&\frac{2}{3}\\
\frac{2}{3}&-\frac{1}{3}&\frac{2}{3}\\
\frac{2}{3}&\frac{2}{3}&-\frac{1}{3}
\end{bmatrix}$
می تواند به عنوان یک انتخاب برای 
$U$
در نظر گرفته شود.
)
  $$\begin{bmatrix}
  -3&1\\
  6&-2\\
  6&-2
  
  \end{bmatrix}$$
 
 \مسئله{} نشان دهید در یک ماتریس مربعی قدر مطلق دترمینان برابر حاصلضرب مقادیر تکین ماتریس است. 
 \\
 \\
 $\star \star$
{\bf  سوالات زیر برای تمرین بیشتر در نظر گرفته شده است و به آن ها نمره ای تعلق نمی گیرد:}
 
 \مسئله{} فرض کنید:
 $$A=\begin{bmatrix}
 0.5&0.2&0.3\\
 0.3&0.8&0.3\\
 0.2&0&0.4
 \end{bmatrix}
 ,v_1=
 \begin{bmatrix}
 0.3\\
 0.6\\
 0.1
 \end{bmatrix}
 ,v_2=
 \begin{bmatrix}
 1\\
 -3\\
 2
 \end{bmatrix}
 ,
 v_3=
 \begin{bmatrix}
 -1\\
 0\\
 1
 \end{bmatrix}
 ,
 w=
 \begin{bmatrix}
 1\\
 1\\
 1
 \end{bmatrix}
 $$
 \begin{enumerate}
 	\item 
 	نشان دهید 
 	$v_1,v_2$
 	و 
 	$v_3$
 	بردار ویژه های 
 	$A$
 	هستند.
 	\item
 	فرض کنید 
 	$x_0$
 	بردار برداری در 
 	$\mathbb{R}^3$
 	باشد که درایه های آن نامنفی باشند و مجموعشان 1 باشد. ثابت کنید وجود دارد ثابت هایی مثل 
 	$c_1,c_2,c_3$
 	که 
 	$x_0=c_1v_1+c_2v_2+c_3v_3$
 	باشد و همچنین 
 	$w^{T}x_0$
 	را محاسبه کنید و نتیجه بگیرید 
 	$c_1=1$.
 	
 	\item 
 	برای 
 	$k=1,2,\cdots$
 	تعریف می کنیم 
 	$x_k=A^kx_0$
 	که 
 	$x_0$
 	در قسمت (2) معرفی شده است .نشان دهید 
 	$x_k\to v_1$
 	زمانی که 
 	$k$
 	افزایش می یابد.
 \end{enumerate}


\مسئله{}فرض کنید 
$A$
یک ماتریس 
$n\times n$
متقارن باشد،فرض 
$x$
هر برداری در 
$\mathbb{C}^n$
باشد و در نظر بگیرید 
$q=\bar{x}^TAx$.
تساوی های زیر نشان می دهند که 
$q=\bar{q}$.
هرکدام از تساوی ها را با ادله کافی توجیه کنید.
$$\bar{q}=\overline{\bar{x}^TAx}=x^T\overline{Ax}=x^TA\bar{x}=(x^TA\bar{x})^T=\bar{x}^TA^Tx=q$$



\مسئله{}
$u\neq 0$
را در 
$\mathbb{R}^n$
در نظر بگیرید،فرض کنید 
$L=span\{u\}$
برای هر 
$y$
در 
$\mathbb{R}^n$
تصویر 
$y$
نسبت به 
$L$
را اینگونه تعریف می کنیم:
$$refl_L \boldsymbol{y}=2.proj_L\boldsymbol{y}-\boldsymbol{y}$$
نشان دهید 
$\boldsymbol{y}\mapsto refl_L\boldsymbol{y}$
یک تبدیل خطی است.


\مسئله{} فرض کنید 
$A=QR$
یک تقسیم بندی 
$QR$
برای ماتریس 
$A$ای
باشد که ستون های آن مستقل خطی هستند.
$A$
را به شکل 
$[A_1 \ \ A_2]$
می نویسیم که 
$A_1$
p ستون دارد.
چگونه می توان یک تقسیم بندی 
$QR$
برای 
$A_1$
یافت؟توضیح دهید تقسیم بندی شما چگونه شرایط یک تقسیم بندی 
$QR$
را حفظ می کند.


\end{document}


